% GENERATED FILE. Edit canonical lessons, then run npm run generate:books.
% canonical-source-hash: fnv1a64:012ede3c8fc24145
% canonical-lessons: ES-C706-asustar, ES-C706-asustado, ES-C706-susto, ES-C706-panico, ES-C706-temer

\chapter{Sustos y miedos}
\label{ch:es-sustos-y-miedos}
\hlchaptermodality{706}

\begin{chapteropening}
\textbf{By the end of this chapter:} \emph{I can tell someone that something gave me a fright or that I fear something.}

Complete the last lesson of chapter 706: I can tell someone that something gave me a fright or that I fear something.
\end{chapteropening}

\section[asustar]{asustar — to frighten, to scare}

\label{lesson:ES-C706-asustar}



\begin{quote}
Before the new one: say the Spanish for ashamed, then the Spanish for to be ashamed.
\end{quote}

\subsection*{You'll want to know: asustar}
\textbf{asustar} — \textquotedblleft{}to frighten, to scare\textquotedblright{}.

\begin{quote}
El perro me asustó.
\end{quote}

\begin{quote}
\emph{The dog frightened me.}
\end{quote}

\emph{Asustarse} is to get scared: \emph{me asusté}.

\begin{cousinweb}
\emph{asustar} is built on \emph{susto} \textquotedblleft{}fright\textquotedblright{} with the prefix \emph{a-}. Where \emph{susto} itself comes from, the origin is not settled; it may be an expressive word.
\end{cousinweb}

\begin{grammarlens}[title={What you've built}]
One more word for telling, explaining and describing at B1.
\end{grammarlens}

\subsection*{Your turn}
\begin{itemize}
\raggedright
  \item \emph{Say it:} \emph{asustar}
  \item \emph{Say it:} \emph{asustar}, once more
  \item \emph{Say it:} a sentence of your own with \emph{asustar}
  \item \emph{Recall:} say the Spanish for ashamed, then the Spanish for to be ashamed, then say \emph{asustar} again
\end{itemize}

\begin{tcolorbox}[breakable,colback=teal!4,colframe=teal!35!black,title={Before you move on}]
What does asustar mean? (\textquotedblleft{}To frighten, to scare\textquotedblright{}.) Say it once more.
\end{tcolorbox}
\section[asustado]{asustado — frightened, scared}

\label{lesson:ES-C706-asustado}



\begin{quote}
Before the new one: say the Spanish for to be ashamed, then the Spanish for to frighten.
\end{quote}

\subsection*{You'll want to know: asustado}
\textbf{asustado} — \textquotedblleft{}frightened, scared\textquotedblright{}.

\begin{quote}
La niña está asustada.
\end{quote}

\begin{quote}
\emph{The little girl is frightened.}
\end{quote}

Used with \emph{estar}.

\begin{cousinweb}
\emph{asustado} is the participle of \emph{asustar}, which is built on \emph{susto} \textquotedblleft{}fright\textquotedblright{}; for \emph{susto} the origin is not settled.
\end{cousinweb}

\begin{grammarlens}[title={What you've built}]
One more word for telling, explaining and describing at B1.
\end{grammarlens}

\subsection*{Your turn}
\begin{itemize}
\raggedright
  \item \emph{Say it:} \emph{asustado}
  \item \emph{Say it:} \emph{asustado}, once more
  \item \emph{Say it:} a sentence of your own with \emph{asustado}
  \item \emph{Recall:} say the Spanish for to be ashamed, then the Spanish for to frighten, then say \emph{asustado} again
\end{itemize}

\begin{tcolorbox}[breakable,colback=teal!4,colframe=teal!35!black,title={Before you move on}]
What does asustado mean? (\textquotedblleft{}Frightened, scared\textquotedblright{}.) Say it once more.
\end{tcolorbox}
\section[el susto]{el susto — the fright, the scare}

\label{lesson:ES-C706-susto}



\begin{quote}
Before the new one: say the Spanish for to frighten, then the Spanish for frightened.
\end{quote}

\subsection*{You'll want to know: el susto}
\textbf{el susto} — \textquotedblleft{}the fright, the scare\textquotedblright{}.

\begin{quote}
¡Qué susto me diste!
\end{quote}

\begin{quote}
\emph{What a fright you gave me!}
\end{quote}

\emph{Dar un susto} — to give someone a fright.

\begin{cousinweb}
For \emph{susto} the origin is not settled. It is old in Spanish and Portuguese, and it may have begun as an expressive word imitating a sharp breath.
\end{cousinweb}

\begin{grammarlens}[title={What you've built}]
One more word for telling, explaining and describing at B1.
\end{grammarlens}

\subsection*{Your turn}
\begin{itemize}
\raggedright
  \item \emph{Say it:} \emph{el susto}
  \item \emph{Say it:} \emph{el susto}, once more
  \item \emph{Say it:} a sentence of your own with \emph{el susto}
  \item \emph{Recall:} say the Spanish for to frighten, then the Spanish for frightened, then say \emph{el susto} again
\end{itemize}

\begin{tcolorbox}[breakable,colback=teal!4,colframe=teal!35!black,title={Before you move on}]
What does el susto mean? (\textquotedblleft{}The fright, the scare\textquotedblright{}.) Say it once more.
\end{tcolorbox}
\section[el pánico]{el pánico — panic}

\label{lesson:ES-C706-panico}



\begin{quote}
Before the new one: say the Spanish for frightened, then the Spanish for the fright.
\end{quote}

\subsection*{You'll want to know: el pánico}
\textbf{el pánico} — \textquotedblleft{}panic\textquotedblright{}.

\begin{quote}
Me da pánico hablar en público.
\end{quote}

\begin{quote}
\emph{Speaking in public terrifies me.}
\end{quote}

\emph{Me da pánico} + infinitive — something terrifies me.

\begin{cousinweb}
\emph{pánico} is usually traced to Greek \emph{panikós} \textquotedblleft{}of Pan\textquotedblright{}, the country god said to cause sudden fear; it came into Spanish through French and Latin.
\end{cousinweb}

\begin{grammarlens}[title={What you've built}]
One more word for telling, explaining and describing at B1.
\end{grammarlens}

\subsection*{Your turn}
\begin{itemize}
\raggedright
  \item \emph{Say it:} \emph{el pánico}
  \item \emph{Say it:} \emph{el pánico}, once more
  \item \emph{Say it:} a sentence of your own with \emph{el pánico}
  \item \emph{Recall:} say the Spanish for frightened, then the Spanish for the fright, then say \emph{el pánico} again
\end{itemize}

\begin{tcolorbox}[breakable,colback=teal!4,colframe=teal!35!black,title={Before you move on}]
What does el pánico mean? (\textquotedblleft{}Panic\textquotedblright{}.) Say it once more.
\end{tcolorbox}
\section[temer]{temer — to fear}

\label{lesson:ES-C706-temer}



\begin{quote}
Before the new one: say the Spanish for the fright, then the Spanish for panic.
\end{quote}

\subsection*{You'll want to know: temer}
\textbf{temer} — \textquotedblleft{}to fear\textquotedblright{}.

\begin{quote}
Temo que llegue tarde.
\end{quote}

\begin{quote}
\emph{I fear he will be late.}
\end{quote}

More formal than \emph{tener miedo}. \emph{Temer que} takes the subjunctive.

\begin{cousinweb}
\emph{temer} is inherited from Latin \emph{timēre} \textquotedblleft{}to fear\textquotedblright{} (EtymDB); English \emph{timid} comes from the same Latin family.
\end{cousinweb}

\begin{grammarlens}[title={What you've built}]
One more word for telling, explaining and describing at B1.
\end{grammarlens}

\subsection*{Your turn}
\begin{itemize}
\raggedright
  \item \emph{Say it:} \emph{temer}
  \item \emph{Say it:} \emph{temer}, once more
  \item \emph{Say it:} a sentence of your own with \emph{temer}
  \item \emph{Recall:} say the Spanish for the fright, then the Spanish for panic, then say \emph{temer} again
\end{itemize}

\begin{tcolorbox}[breakable,colback=teal!4,colframe=teal!35!black,title={Before you move on}]
What does temer mean? (\textquotedblleft{}To fear\textquotedblright{}.) Say it once more.
\end{tcolorbox}
